\section{Experiments}
\label{sec:exp}

We organize experiments around six questions. \textbf{RQ1}: does user-conditioned bilinear reranking beat a user-agnostic baseline on session-level ranking? \textbf{RQ2}: is the bilinear term $h^\top A u$ necessary, or does a user-only bias $w^\top u$ suffice? \textbf{RQ3}: does the construction of $u$ matter --- does an LLM-extracted typed schema beat behavioral features or random vectors? \textbf{RQ4}: is the gain stable across random seeds? \textbf{RQ5}: is the bilinear form specifically necessary, or would any coupling of $h$ and $u$ (concatenation, FiLM) work? \textbf{RQ6}: how sensitive is the gain to the user-vector dimension $u_{\dim}$?

\subsection{Setup}

\paragraph{Dataset.}
We use AOL-IA~\cite{pass2006_picture} via \texttt{ir\_datasets}, sampling 200k sessions (each with at least four query events) from 36M queries across 657K users. Table~\ref{tab:funnel} reports the full data funnel. The evaluation set of $n{=}201$ test sessions is obtained after two filters at candidate-set construction: (i) the eval query must have at least nine other documents clicked by different users in the same split (to form the negative pool), and (ii) all ten candidates must have complete precomputed backbone features. Long-tail queries with few cross-session clicks fail the first filter; this is a known limitation of click-log evaluation without full SERP impressions, which AOL-IA does not provide.

\begin{table}[t]
\centering
\caption{Data funnel from raw AOL-IA to the final evaluation set. The test-set shrinkage is a property of click-log evaluation, not of the method; we report it transparently.}
\label{tab:funnel}
\begin{tabular}{lr}
\toprule
Stage & Count \\
\midrule
AOL-IA raw queries & 36{,}389{,}567 \\
Anonymized users & 657{,}427 \\
Sessions ($\geq$4 events, sampled) & 200{,}000 \\
Train pairs (position-debiased) & 3{,}305 \\
Dev pairs & 344 \\
Test pairs & 353 \\
Eval sessions (after candidate filter) & 201 \\
\bottomrule
\end{tabular}
\end{table}

\paragraph{Metrics.}
Session-level nDCG@10 is the primary metric (per-session nDCG averaged across sessions), with MRR and Recall@10 as secondary. We report paired bootstrap 95\% confidence intervals on the nDCG@10 difference between each system and the user-agnostic baseline. The bootstrap resamples sessions with replacement (10{,}000 replicates), recomputeing the mean per-session nDCG@10 difference each time; pairing is by session, so each system and the baseline are scored on identical candidate sets.

\paragraph{Candidate set construction.}
For each test query, the candidate set consists of the single clicked document (relevance 1) and nine documents clicked for the same query by other users but not by this user (relevance 0). This follows the click-log convention; we disclose in Sec.~\ref{sec:discussion} that these negatives are cross-session clicked distractors, a proxy for true non-relevance. A document clicked by another user for the same query may in fact be relevant to this user too; we treat this as a known noise source and discuss sensitivity to negative-sampling choices as future work.

\paragraph{Systems.}
\textbf{Baseline} is the frozen BM25+BGE query-flow scorer with $s_{\mathrm{base}}$ only. \textbf{Pref-Rerank} variants: \texttt{full} ($s_{\mathrm{base}} + w^\top u + h^\top A u$), \texttt{no\_user\_bias} ($s_{\mathrm{base}} + h^\top A u$), \texttt{linear\_only} ($s_{\mathrm{base}} + w^\top u$), all with the LLM-extracted typed schema as $u$. Representation ablations use the \texttt{no\_user\_bias} score with $u$ replaced by \texttt{classic} (behavioral features) or \texttt{random\_schema} (random vector).

\subsection{RQ1: Main Result}
\label{sec:exp:main}

Table~\ref{tab:main} reports the main comparison. The user-agnostic baseline reaches nDCG@10 $= 0.459$; Pref-Rerank with behavioral features (\texttt{classic}) reaches $0.543$, a gain of $+0.084$ with paired bootstrap 95\% CI $[+0.034, +0.131]$, which excludes zero. This supports \textbf{C1}: a 12{,}288-parameter bilinear intervention on a frozen scorer produces a statistically significant improvement in session-level evidence ranking at $n{=}201$ test sessions.

\begin{table}[t]
\centering
\caption{Main result on 200k AOL sessions ($n{=}201$ test sessions). Pref-Rerank (\texttt{classic}: behavioral features + bilinear $h^\top A u$) vs.\ the user-agnostic frozen BM25+BGE baseline. $\Delta$ is the per-session nDCG@10 difference; CI from paired bootstrap (10k resamples). MRR reported for completeness; Recall@10 is 1.0 for all systems by construction (single relevant doc in a pool of 10).}
\label{tab:main}
\begin{tabular}{lrrr}
\toprule
System & nDCG@10 & MRR & $\Delta$ nDCG@10 (95\% CI) \\
\midrule
Baseline (user-agnostic) & 0.459 & 0.304 & --- \\
\textbf{Pref-Rerank (\texttt{classic}, ours)} & \textbf{0.543} & \textbf{0.408} & \textbf{+0.084 [+0.034, +0.131]} \\
\bottomrule
\end{tabular}
\end{table}


\subsection{RQ2: Mechanism --- the bilinear term is the active component}
\label{sec:exp:mechanism}

Table~\ref{tab:mechanism} isolates which score term drives the gain. \texttt{linear\_only} (user-only bias $w^\top u$, no bilinear) is indistinguishable from baseline ($\Delta = 0.000$), confirming that a document-independent user offset contributes nothing for ranking, which is comparative within a candidate set. The \texttt{no\_user\_bias} variant (bilinear only) and the \texttt{full} variant (both terms) both reach $0.513$, showing that adding $w^\top u$ contributes nothing on top of the bilinear term. This supports \textbf{C2}: the bilinear term $h^\top A u$ is the active mechanism; the user-only bias $w^\top u$ is redundant and can be dropped, yielding the 12{,}288-parameter form.

\begin{table}[t]
\centering
\caption{Mechanism ablation (RQ2): which score term drives the gain? Only the bilinear term $h^\top A u$ contributes; the user-only bias $w^\top u$ alone is indistinguishable from baseline, and adding it on top of the bilinear term yields no further gain. $\Delta$ vs.\ user-agnostic baseline; CI from paired bootstrap.}
\label{tab:mechanism}
\begin{tabular}{lcrr}
\toprule
Variant & Score formula & nDCG@10 & $\Delta$ (95\% CI) \\
\midrule
Baseline & $s_{\mathrm{base}}$ & 0.459 & --- \\
\texttt{linear\_only} & $s_{\mathrm{base}} + w^\top u$ & 0.459 & +0.000 [+0.000, +0.000] \\
\texttt{no\_user\_bias} & $s_{\mathrm{base}} + h^\top A u$ & 0.513 & +0.054 [+0.013, +0.094] \\
\texttt{full} & $s_{\mathrm{base}} + w^\top u + h^\top A u$ & 0.513 & +0.054 [+0.013, +0.094] \\
\bottomrule
\end{tabular}
\end{table}


\subsection{RQ3: Which user representation works?}
\label{sec:exp:representation}

Table~\ref{tab:representation} compares three constructions of $u$ under the \texttt{no\_user\_bias} score. Behavioral features (\texttt{classic}) yield the strongest gain ($+0.084$, CI $[+0.034, +0.131]$). The LLM-extracted typed schema (\texttt{no\_user\_bias}) is also significant ($+0.054$, CI $[+0.013, +0.094]$) but weaker. Crucially, a random user vector (\texttt{random\_schema}) is \emph{not} significant ($+0.027$, CI $[-0.010, +0.065]$, crosses zero).

\begin{table}[t]
\centering
\caption{User-representation ablation (RQ3): all three constructions of $u$ use the \texttt{no\_user\_bias} bilinear score. Behavioral features yield the strongest and tightest gain; the LLM-extracted typed schema is also significant but weaker; a random vector is not significant (CI crosses zero), ruling out pure session-identity memorization.}
\label{tab:representation}
\begin{tabular}{llrr}
\toprule
$u$ construction & Dim & nDCG@10 & $\Delta$ (95\% CI) \\
\midrule
Baseline & --- & 0.459 & --- \\
Random vector (\texttt{random\_schema}) & 16 & 0.486 & +0.027 [$-$0.010, +0.065] \\
LLM typed schema (\texttt{no\_user\_bias}) & 16 & 0.513 & +0.054 [+0.013, +0.094] \\
\textbf{Behavioral features (\texttt{classic})} & 6 $\to$ 16 & \textbf{0.543} & \textbf{+0.084 [+0.034, +0.131]} \\
\bottomrule
\end{tabular}
\end{table}


This pattern addresses the central confound head-on. If the bilinear term merely memorized session identity, random user vectors should work just as well as real ones. They do not. The bilinear interaction exploits genuine user signal: behavioral features, which carry real session statistics, outperform both the typed schema and random vectors. The LLM schema retains an interpretability advantage that nDCG does not measure, exposing a human-readable four-dimensional preference (e.g., ``official + detailed + data + latest''), but it does not provide a ranking-quality edge over simpler behavioral features.

\subsection{RQ4: Stability across seeds}
\label{sec:exp:seed}

A natural concern is whether the $+0.084$ gain is an artifact of a single initialization. Table~\ref{tab:seed} reports the \texttt{classic} variant under three independent seeds. The gain is stable: $\Delta \in \{+0.072, +0.084, +0.087\}$, mean $+0.081$ with standard deviation $0.008$, and every seed's paired bootstrap 95\% CI excludes zero. The improvement is not sensitive to the random seed.

\begin{table}[t]
\centering
\caption{Multi-seed stability (RQ4). The \texttt{classic} variant trained under three independent seeds. The gain is stable: mean $\Delta{=}{+}0.081$, std $0.008$, and every seed's 95\% CI excludes zero. The result is not an artifact of a single initialization.}
\label{tab:seed}
\begin{tabular}{lrrr}
\toprule
Seed & nDCG@10 & $\Delta$ & 95\% CI \\
\midrule
42 & 0.543 & +0.084 & [+0.034, +0.131] \\
1  & 0.531 & +0.072 & [+0.028, +0.117] \\
2  & 0.546 & +0.087 & [+0.039, +0.137] \\
\midrule
\textbf{mean $\pm$ std} & \textbf{0.540 $\pm$ 0.008} & \textbf{+0.081 $\pm$ 0.008} & all CIs exclude zero \\
\bottomrule
\end{tabular}
\end{table}


\subsection{RQ5: Is the bilinear form necessary, or would any $h$--$u$ coupling do?}
\label{sec:exp:mech}

The bilinear term $h^\top A u$ is one of several ways to couple the document representation $h$ with a user vector $u$. Table~\ref{tab:pareto} compares it against two equally minimal alternatives under the \texttt{classic} user representation: concatenation-plus-linear ($\mathrm{Linear}([h;u]) \to 1$, 785 parameters) and FiLM-style modulation ($\gamma,\beta = Wu$ applied elementwise to $h$, 26{,}112 parameters). Concatenation is significant but weaker ($+0.054$); FiLM matches the bilinear gain ($+0.081$) at $2.1\times$ the parameter count. The bilinear form is Pareto-optimal on this benchmark: it reaches the strongest gain with fewer parameters than FiLM and a larger gain than concatenation.

\begin{table}[t]
\centering
\caption{Mechanism comparison (RQ5): the bilinear form against two equally minimal alternatives that also couple $h$ and $u$, under the \texttt{classic} user representation. Concatenation-plus-linear (785 params) is significant but weaker; FiLM-style modulation (26{,}112 params) matches the bilinear gain at $2.1\times$ the parameters. The bilinear form is Pareto-optimal: fewest parameters, strongest gain.}
\label{tab:pareto}
\begin{tabular}{lrrrr}
\toprule
Mechanism & Params & nDCG@10 & $\Delta$ & 95\% CI \\
\midrule
Baseline & 0 & 0.459 & --- & --- \\
concat + linear & 785 & 0.513 & +0.054 & [+0.011, +0.095] \\
FiLM ($\gamma,\beta = Wu$) & 26{,}112 & 0.540 & +0.081 & [+0.031, +0.128] \\
\textbf{bilinear $h^\top A u$ (ours)} & \textbf{12{,}304} & \textbf{0.543} & \textbf{+0.084} & \textbf{[+0.034, +0.131]} \\
\bottomrule
\end{tabular}
\end{table}


\subsection{RQ6: How sensitive is the gain to $u_{\dim}$?}
\label{sec:exp:dim}

Table~\ref{tab:dim} sweeps the user-vector dimension $u_{\dim} \in \{4,8,16,32,64\}$, zero-padding the six behavioral features and resizing $A$ accordingly. The gain saturates at $u_{\dim}{=}8$: dimensions beyond that do not improve nDCG and only inflate the trainable budget (from $6{,}152$ to $49{,}216$ parameters). Our default $u_{\dim}{=}16$ sits safely on the plateau, confirming it was not tuned to a lucky value.

\begin{table}[t]
\centering
\caption{User-vector dimension ablation (RQ6), \texttt{classic} variant. The 6 behavioral features are zero-padded to $u_{\dim}$ and the bilinear matrix $A$ resized accordingly. The gain saturates at $u_{\dim}{=}8$: further capacity does not help and only inflates the parameter count. $u_{\dim}{=}16$ (our default) sits safely on the plateau.}
\label{tab:dim}
\begin{tabular}{rrrrr}
\toprule
$u_{\dim}$ & Params & nDCG@10 & $\Delta$ & 95\% CI \\
\midrule
4  & 3{,}076  & 0.525 & +0.066 & [+0.019, +0.112] \\
8  & 6{,}152  & 0.543 & +0.084 & [+0.034, +0.131] \\
16 & 12{,}304 & 0.543 & +0.084 & [+0.034, +0.131] \\
32 & 24{,}608 & 0.543 & +0.084 & [+0.034, +0.131] \\
64 & 49{,}216 & 0.543 & +0.084 & [+0.035, +0.132] \\
\bottomrule
\end{tabular}
\end{table}

